\section{Deliverables}

\begin{enumerate}
\item Write a Python module, called \texttt{feedback\_supression.py},
  that inherits from \texttt{buffer.py} and that implements one of the
  proposed solutions (or any other you prefer). Define a new
  command-line parameter to switch on this functionality, that by
  default should be disabled (off). If necessary, use extra parameters
  that can be controlled in real time.
\item Insert the code of \texttt{feedback\_supression.py} in a Jupyter
  notebook cell and generate the archive
  \texttt{feedback\_supression.py} using \texttt{\%\%writefile
    feedback\_supression.py}. The notebook should also include the
  names of the members of the working group and I example of how to
  use your implementation.
\item Notice that for testing your implementation it is essential that
  your operating system does not cancel the echo. If your OS cancel
  the echo, you must perform the tests in Linux (see
  \href{https://tecnologias-multimedia.github.io/contents/framework/}{Framework}).
\end{enumerate}

\begin{comment}
More concretely:

\begin{enumerate}

\item If you implement the ``$s(t)$ delay and subtract solution'',
  you must estimate $a$, $d$ and perform the signals subtraction to
  remove the feedback signal ($s(t)$). 

\item In a ``$s(t)$ delay and subtract solution'', if you also
  consider the frequency response of the near-end audioset to estimate
  a better feedback signal, first you will need to find ${\mathbf H}$ (the
  discrete frequency response of your audioset). For this, the
  near-end speaker should generate an impulse signal
  ${\mathbf \delta}$, and in absence of any other sound signal, record
  the echo and compute its Fourier transform (it is a good idea to
  repeat this process several times to obtain a better estimation of
  ${\mathbf H}$). Finally, .

  Notice that the frequency characterization of the far-end audioset
  can be also used in a ``$n(t)$ delay and subtract
  solution''. Remember that filtering operations must be implemented
  with convolutions in the temporal domain, but are dot products in
  the frequency domain.
  
\item Use LMS to find a estimation of the echo signal and perform the
  echo cancellation. In this case, consider the use of a thread to
  compute the coefficients of the LMS filter
  (Equation~\eqref{eq:update} can run with a cadence much higher than
  a sample-time) and compute the estimated feedback signal
  (Equation~\eqref{eq:LMS_feedback}) for every chunk. Notice that for
  doing that, we will require the first $k$ samples of the next chunk.

\item 
\end{enumerate}

Optionally, but interesting for your mark, use any other technique
(for example, an artificial neural network\footnote{ADALINE is the
  simplest
  \href{https://en.wikipedia.org/wiki/Neural_network_(machine_learning)}{AAN}
  ever developed!}) for estimating the echo, and use it for removing
the echo (obviously, in real-time). Take also into consideration that
the parameters that determine the estimation of the echo signal should
be continously\footnote{A 1-seconds cadence should be enought.}
monitored becase the physical composition of the audiosets can
be dynamic (for example, the inclination of the screen of our laptop
can be modified at any moment).

Finally, notice that
operation can help to fine-tune $d$.

% Mark: \textbf{10 points}.

\item To the report describing how your module works, \textbf{add a
    section evaluating it} (with a ratio out of 10), considering that
  each of the following rubrics has the same weight in the evaluation
  score:
  \begin{enumerate}
  \item Compared, for example, with the feedback supression provided
    by Google Meet, how of effective is your proposal? You must
    describe how this comparison has been carried out.
  \item  Is your implementation efficient in terms of CPU usage? Explain.
  \item Is the code of the module properly explained and
    documented (in the own module)?
  \end{enumerate}
\end{enumerate}
\end{comment}
